\chapter{Discussion}\label{ch:discussion}

The core chapters each answered a focused question with a specific method, benchmark and set of experiments. This chapter steps back and interprets their findings together. Section~\ref{sec:disc:rq} answers the four research questions of Chapter~\ref{ch:intro}; the last of them, RQ4, is answered as a synthesis and research agenda whose elements are graded by the strength of their evidence. Section~\ref{sec:disc:method} reflects on the research framework of Chapter~\ref{ch:methodology}: what it delivered and where it fell short. Section~\ref{sec:disc:limits} states the limitations of the thesis, and Section~\ref{sec:disc:ethics} discusses its broader implications.

% ==============================================================================
\section{Answering the Research Questions}\label{sec:disc:rq}

\subsection{RQ1: how should a physical agent represent the 3D world?}

Chapter~\ref{ch:occnet} argued that a physical agent should describe its surroundings by 3D semantic occupancy rather than by lists of boxes. The evidence supports this in three respects, each with a qualification. First, occupancy describes shapes that boxes cannot: irregularly shaped vehicles are segmented better when supervised with occupancy than with voxelised boxes, and the gap widens as resolution increases. These objects, however, belong to classes seen in training, so the evidence concerns shape, not novelty. Second, occupancy works as an interface to other tasks: the same representation improves semantic scene completion, supports LiDAR segmentation from cameras, serves as a pre-training target for detection, and lowers the collision rate of a planner. The planning evidence is the weakest of these: it comes from open-loop evaluation on nuScenes with the ST-P3 protocol, the differences in collision rate are small in absolute terms (0.06 to 0.24 percentage points relative to planning on boxes), and each configuration was trained once. The benefit also comes from occupancy as the planner's \emph{input}: features pre-trained on occupancy did not improve planning when fine-tuned. Third, the representation could be scaled: the OpenScene redistribution of nuPlan turned occupancy from a 5.5-hour academic benchmark into more than 120 hours of supervision that the community has used in end-to-end driving and world-model challenges.

The answer to RQ1 is therefore that occupancy makes the representation \emph{geometrically open}, in the sense that it can express the shape of whatever is in the scene, and \emph{queryable by downstream tasks}, and that it can be supervised at scale only through a deliberate investment in data. It does not make the representation semantically open: OpenOcc labels voxels with a fixed set of sixteen classes, so an unseen object is occupied space without a meaning. Robustness of occupancy perception to distribution shift, for example across cities or sensor set-ups, was not evaluated in this thesis.

\subsection{RQ2: can structured reasoning with vision-language models help a driving agent generalise?}

Chapter~\ref{ch:drivelm} gave a qualified yes, and separating two effects is essential to reading it correctly. The first effect is that of the pre-trained vision-language backbone. In zero-shot transfer from nuScenes to the Waymo sensor configuration, the displacement error of the vision-language baseline was 2.78 without any question answering and 2.76 without reasoning context, against 4.16 for the driving-specific UniAD-Single; most of the transfer advantage over UniAD-Single therefore comes from the backbone. The second effect is that of graph-structured reasoning, which is smaller: with graph context the displacement error fell further to 2.63, with most of the gain in the accuracy of speed decisions (from 43.90 to 54.29). In-domain on nuScenes, reasoning did not help: the variant without context predicted behaviour most accurately (61.45 against 57.49 for graph context) and with the lowest displacement error (1.39 against 1.74). On CARLA scenes with pedestrians unseen during training, behaviour accuracy collapsed (from 59.63 to 4.59) and recovered to 27.04 only when a question specifically about pedestrians was added to the graph, at the cost of accuracy on regular scenes (from 59.63 to 52.17). All of these are single runs.

The pattern is itself a finding: explicit reasoning helps in the long tail rather than in the average case, and it helps most when the reasoning is directed at the right entity, which in this experiment was chosen by someone who knew the shift. It also exposes two costs. The first is latency: the reasoning agent ran at 0.16 frames per second, far too slow for control; the co-authored ETA later addressed this with a fast-and-slow design in which a large model reasons one step ahead while a small model reacts~\cite{hamdan2025eta}. The second is reliability: DriveBench showed that vision-language models can retain most of their score with the image removed, and that 78.6\% of DriveLM behaviour answers are ``Going Straight''~\cite{xie2025vlms}. The answer to RQ2 is therefore that structured reasoning adds measurable gains under shift on top of what a pre-trained backbone provides, provided it is grounded, evaluated on balanced data, and directed at the relevant entities; it is not a substitute for broad pre-training, and it brings no benefit in familiar situations.

\subsection{RQ3: how can a learned policy stay reliable under shift?}

For driving (RQ3a), Chapter~\ref{ch:centaur} showed that a planner can improve itself at test time without rules or labels, in non-reactive simulation. Minimising Cluster Entropy with a single test-time gradient step raised the NAVSIM driving score of Hydra-MDP from 90.3 to 92.6 in a single run, and the three-seed mean of 92.10 $\pm$ 0.33 ranked first among three-seed entries on the NAVSIM v1.1 leaderboard at the time of submission. Two qualifications matter. The margin over the semantic-entropy variant of the same approach is small (0.23 PDMS over three seeds) and was not tested for significance; and every test-time training variant in the chapter improved its base planner, by 0.8 to 2.3 PDMS, whereas every fallback variant lowered it. The more robust conclusion is therefore about the mechanism rather than the measure: learned test-time adaptation on the planner's own uncertainty improves safety sub-scores with little loss of progress (ego progress changed from 86.5 to 85.9), whereas hand-designed guards buy safety by giving up progress (the occupancy-based optimisation of UniAD lowers its score from 83.4 to 75.2, and fallback layers reduce a strong planner to about 63--66). Test-time adaptation does not close the gap on the hardest cases: on the safety-critical \texttt{navsafe} split, the adapted planner reaches 74.14 against 94.65 for human driving, and it performs worse than its base planner on the yielding category (59.3 against 71.6).

For manipulation (RQ3b), Chapter~\ref{ch:kai0} located the problem not only in the amount of data or computation but in inconsistency between the distribution of demonstrations, what the policy can represent, and what is executed on hardware. Addressing these inconsistencies with Model Arithmetic, Stage Advantage and Train--Deploy Alignment, on top of the heavily pre-trained $\pi_{0.5}$ policy and with about 20 hours of data per task, raised the success rate on the T-shirt folding task from about 27\% to about 97\% over 30 real-robot trials, with a roughly tenfold increase in throughput (values read from the source's figure). The co-authored AgiBot World reports, in contrast, that pre-training on more and verified robot data improves both in-distribution and out-of-distribution success~\cite{bu2025agibot_iros}. The two findings are better read as complementary than as contradictory. Scale buys \emph{coverage} of the deployment distribution, whereas alignment buys \emph{consistency} between how a policy is trained and how it is used; the same AgiBot World study also found that its 528 verified trajectories alone outperformed all 1,010 trajectories including 482 unverified ones (0.59 against 0.41). The answer to RQ3 is therefore that, in the settings studied, reliability under shift improved when the loop between training and deployment was closed, either online, by adapting on the agent's own uncertainty, or offline, by aligning data, model and execution. Rules reduced progress in the driving experiments, and alignment improved robustness at a modest post-training budget in the manipulation experiments; the trade-off between alignment and further scaling was not tested.

\subsection{RQ4: synthesis and research agenda for self-correcting agents}

RQ4 asks what design lessons for self-correcting physical agents follow from the answers to RQ1--RQ3. The lessons can be assembled into the anatomy sketched in Fig.~\ref{fig:disc:anatomy}. It is a synthesis, not a demonstrated system: each element is motivated by the findings of this thesis and by related work, but none has been evaluated as part of an integrated agent. Table~\ref{tab:disc:evidence} grades the evidence for each element.

\begin{figure}[t]
  \centering
  \tikzfig{ch-discussion/figures/self_correcting}
  \caption[Anatomy of a self-correcting physical agent]{Anatomy of a self-correcting physical agent, proposed as a synthesis of this thesis. A universal representation feeds a fast, planning-oriented policy directly and a slow, explicit reasoner selectively. A competence monitor computes signals of the agent's reliability and uses them to adapt the policy at test time, to invoke slow reasoning, or to request more perception. A learning loop feeds residual failures back into data, alignment and benchmarks. The evidence for each element is graded in Table~\ref{tab:disc:evidence}; works marked $^{*}$ are co-authored rather than led by the author.}
  \label{fig:disc:anatomy}
\end{figure}

\begin{table}[t]
  \centering
  \caption[Evidence for the elements of the proposed self-correcting agent]{Evidence for the elements of the proposed self-correcting agent (Fig.~\ref{fig:disc:anatomy}). Source: \emph{led}, work led by the author and presented in this thesis; \emph{co-authored}, co-authored work not presented as a contribution of this thesis; \emph{literature}, work by others. Status: \emph{demonstrated} in the setting stated; \emph{partial}, some components demonstrated; \emph{hypothesised}, not tested.}
  \label{tab:disc:evidence}
  \footnotesize
  \renewcommand{\arraystretch}{1.2}
  \begin{tabularx}{\textwidth}{@{}>{\raggedright\arraybackslash}p{2.3cm} >{\raggedright\arraybackslash}X >{\raggedright\arraybackslash}p{1.75cm} >{\raggedright\arraybackslash}p{1.25cm} >{\raggedright\arraybackslash}p{2.35cm}@{}}
    \toprule
    Element & Evidence & Source & Fidelity & Status \\
    \midrule
    Open, queryable representation & occupancy expresses irregular shapes and helps planning (Ch.~\ref{ch:occnet}); planning-oriented design (UniAD) & led; co-authored & L0 & partial (geometry only) \\
    Fast and slow decisions & reasoning helps under shift but is slow (Ch.~\ref{ch:drivelm}); asynchronous large and small models (ETA) & led; co-authored & L0; L2 (ETA) & partial \\
    Competence monitor & label-free behaviour-mode uncertainty (Cluster Entropy, Ch.~\ref{ch:centaur}); supervised stage progress, used in training (Stage Advantage, Ch.~\ref{ch:kai0}) & led & L1; L3 & partial (no monitor that triggers reasoning) \\
    Adaptation at deployment & test-time training of the planner (Ch.~\ref{ch:centaur}) & led & L1 & demonstrated in non-reactive simulation \\
    Learning loop & recovery data and train--deploy alignment (Ch.~\ref{ch:kai0}); benchmark iterations (Ch.~\ref{ch:methodology}); self-directed learning (perspective) & led; co-authored & L3; -- & partial; agent-level loop hypothesised \\
    \bottomrule
  \end{tabularx}
\end{table}

\paragraph{An open, queryable representation.} Downstream modules should read the scene from a representation that does not presuppose which objects matter (Chapter~\ref{ch:occnet}), and they should be trained towards the final task so that upstream errors can be absorbed downstream, as the co-authored UniAD demonstrated~\cite{hu2023uniad}. The evidence is limited to geometry and to open-loop planning.

\paragraph{Two speeds of decision.} Explicit reasoning helps under shift but is slow (Chapter~\ref{ch:drivelm}). An agent could therefore pair a fast policy, which acts on every step, with a slow reasoner that is consulted when the situation is unfamiliar or that thinks ahead asynchronously. This element rests mainly on co-authored work~\cite{hamdan2025eta}; no chapter of this thesis builds such a system.

\paragraph{A competence monitor.} The agent needs a signal of its own reliability. Chapter~\ref{ch:centaur} provides a label-free one: disagreement among the planner's behaviour modes, measured by Cluster Entropy, which also flags a share of failures, though with an accuracy below that of a trivial detector because failures are rare. Chapter~\ref{ch:kai0} provides a different kind: task progress by stage, which is supervised (it requires manually annotated stage labels) and is used during training rather than at deployment. A monitor that decides online when to adapt, when to reason slowly and when to gather more information remains hypothetical.

\paragraph{Adaptation at deployment.} The monitor must be connected to a mechanism that changes behaviour. Chapter~\ref{ch:centaur} demonstrates one, a test-time update of the policy, in non-reactive simulation; hand-designed fallbacks are the degenerate case of this mechanism, and the evidence of that chapter suggests that learned adaptation dominates them. \chiz does not adapt at deployment; its robustness comes from training-side alignment.

\paragraph{A learning loop that outlives deployment.} Residual failures should flow back into data, into the alignment of training with execution, and into new benchmarks. At the level of the research community, this loop is the cycle of Chapter~\ref{ch:methodology}; at the level of a single agent, it would require the agent itself to collect and learn from its failures. The co-authored perspective of~\cite{chen2025robot} describes the agent-level version: a robot that identifies its own goals, acquires skills to reach them, and monitors progress with a learned value function. Neither Centaur nor \chiz reaches that point, since both depend on demonstrations, and \chiz on human-triggered recovery data; each implements a piece of it.

The answer to RQ4 is therefore a research agenda rather than a system: represent openly, decide at two speeds, monitor competence, adapt at deployment, and keep learning from residual failures. Two of these elements, test-time adaptation and training-side alignment, are demonstrated in this thesis in the settings stated; the others are partially supported or remain hypotheses to be tested.

% ==============================================================================
\section{Reflections on the Research Framework}\label{sec:disc:method}

The benchmark-driven cycle of Chapter~\ref{ch:methodology} delivered what it promised in one respect: it multiplied the effort spent on each problem beyond what one group could provide. The ``Driving with Language'' challenge attracted 152 teams, OpenScene became the data source of later benchmarks and challenges, and the lane, topology and occupancy benchmarks were each extended by a later project. Three lessons, however, qualify this success.

\paragraph{Benchmarks carry the biases of their construction.} DriveBench's audit of DriveLM-Data found strong answer imbalance and showed that models could score without looking at the image~\cite{xie2025vlms}; open-loop planning metrics on nuScenes were shown by others to be solvable from the ego vehicle's own state~\cite{zhai2023rethinking,li2024ego}. A benchmark built to expose one shift can hide another. In hindsight, balanced answer distributions, input-ablation baselines (``what does a model score without the image?'') and dedicated shift splits should be part of a benchmark's first release rather than of a later audit; the \texttt{navsafe} split of Chapter~\ref{ch:centaur} is an example of building the stress test in from the start.

\paragraph{Fidelity must be matched to the claim, and stated.} Claims about perception and reasoning were made at the open-loop level, claims about driving policies in non-reactive simulation, and only claims about manipulation on real hardware. This matching was deliberate, but it means that none of the driving results of this thesis has been validated in reactive closed loop or on a vehicle. The ladder of Fig.~\ref{fig:method:ladder} is useful precisely because it makes this gap visible.

\paragraph{Scale and consistency are both resources.} The research began with a strong emphasis on scale: larger benchmarks, larger models, larger challenges. Its final chapter found that, for real-world manipulation, the consistency of data, model and execution mattered as well as their amount. A research framework for physical agents should therefore treat data quality and alignment as objects of measurement, alongside dataset size.

% ==============================================================================
\section{Limitations}\label{sec:disc:limits}

The conclusions of this thesis are subject to the following limitations, in addition to the threats to validity of Section~\ref{sec:method:validity}.

\begin{itemize}[leftmargin=*, itemsep=3pt]
  \item \textbf{Domains.} The thesis studies two kinds of physical agents, road vehicles and bimanual manipulators, and within manipulation only three garment-handling tasks on two robot platforms. Other embodiments may face shifts not covered by the taxonomy of Chapter~\ref{ch:methodology}.
  \item \textbf{Evaluation fidelity for driving.} No learned driving agent was evaluated in reactive closed loop or on a vehicle. DriveLM was evaluated open-loop, and NAVSIM approximates closed-loop behaviour with non-reactive rollouts.
  \item \textbf{Shift coverage.} Chapter~\ref{ch:occnet} does not test distribution shift, and the unseen-object experiment of Chapter~\ref{ch:drivelm} relies on a question written with knowledge of the shift.
  \item \textbf{Statistics.} Most results are single runs without significance tests, and several margins (for example between Cluster Entropy and semantic entropy) are within plausible run-to-run variation. The real-robot results of Chapter~\ref{ch:kai0} are available only as bar charts in the source; the system-level comparison on Task~A is tabulated from values read from the chart, accurate to about one unit.
  \item \textbf{Semantic openness.} Occupancy removes the closed-world assumption from geometry but not from semantics, and the reasoning of DriveLM relies on a vision-language model whose knowledge and failure modes are inherited from its pre-training data.
  \item \textbf{Guarantees.} Test-time training and train--deploy alignment improve average behaviour and robustness, but they provide no formal safety guarantee, and an adaptation step can make the policy worse in some categories, as the yielding results on \texttt{navsafe} show.
  \item \textbf{Resources and dependence on pre-training.} Several results depend on large pre-trained models and on multi-GPU training, which limits their reproducibility for groups with fewer resources; the budget of Chapter~\ref{ch:kai0} is a post-training budget on top of a heavily pre-trained policy.
\end{itemize}

% ==============================================================================
\section{Broader Implications}\label{sec:disc:ethics}

Robust physical agents promise safer roads and useful robots, but they also concentrate responsibility in systems whose behaviour under rare conditions is hard to predict. Several implications of this thesis bear on that tension. First, explainability is not the same as reliability: an agent that explains its decisions in fluent language may still be answering from priors rather than from what it sees, so explanations should be audited as rigorously as actions. Second, self-correction at deployment changes the object that is certified: a policy that adapts online is no longer the policy that was tested. The regression of the adapted planner on yielding in \texttt{navsafe} is a reminder that such changes can affect other road users, such as pedestrians, in ways that aggregate scores hide. Regulation and testing practice will need to account for adaptive behaviour, for instance by bounding the size of test-time updates, by evaluating adapted policies per scenario category, or by monitoring competence signals such as those studied here. Third, the data behind these systems depend on human work: annotators who verify labels and write questions, and teleoperators who demonstrate and correct robot behaviour; their contribution, and their working conditions, deserve recognition in how such data are produced and credited. Finally, the benchmarks released in this thesis are derived from public datasets collected in a small number of cities; performance on them should not be taken as evidence of performance in regions or conditions they do not represent.
